% =====================================================================
% Sekcja 1 - Wprowadzenie
% =====================================================================

\section{Wprowadzenie}
\label{sec:wprowadzenie}

Perplexity (PPL) powstało jako miara jakości modelu językowego -
mówi, jak bardzo model jest zaskoczony obserwowanym tekstem. Z
czasem zaczęto je jednak traktować szerzej: jako tani sygnał o tym,
czy model dobrze radzi sobie z konkretnym zadaniem, liczony bez
potrzeby sprawdzania wyniku innym modelem czy człowiekiem. Pomysł
jest kuszący zwłaszcza w agentach LLM, gdzie taki sygnał mógłby
sterować decyzjami - kazać agentowi wstrzymać się z odpowiedzią,
poprosić o pomoc albo przeplanować działanie, gdy jego pewność spada.

Literatura nie jest tu jednak zgodna. Część prac pokazuje, że niższe
PPL wiąże się z lepszymi wynikami, inne - że PPL bywa zawodne i
podatne na artefakty długości czy interpunkcji. Większość tych badań
dotyczy przy tym pojedynczych, statycznych promptów. Tymczasem
agent nie produkuje jednej odpowiedzi, lecz całą trajektorię -
naprzemienne kroki rozumowania, wywołań narzędzi i obserwacji - i nie
jest oczywiste, czy wnioski z analizy pojedynczych promptów dają się
na taki przypadek przenieść.

Niniejsza praca bada to empirycznie. Sprawdzane jest, czy perplexity
i pokrewne sygnały rozkładowe (entropia tokenowa, top-1 margin)
przewidują poprawność odpowiedzi małego modelu LLM na zadaniach
matematycznych z GSM8K, w dwóch ustawieniach: czystego
chain-of-thought oraz agenta ReAct z kalkulatorem. Osobno badane
jest, w którym miejscu trajektorii agenta sygnał ten jest
najsilniejszy.

Z eksperymentów wyłaniają się trzy obserwacje. Metryki korzystające
z całego rozkładu prawdopodobieństw (entropia, margin) konsekwentnie
przewidują poprawność lepiej niż samo perplexity. Wewnątrz agenta
sygnał nie jest rozłożony równomiernie - najsilniej koreluje z
poprawnością pewność na segmentach formułujących wywołania narzędzia,
a praktycznie zanika na finalnej odpowiedzi. Wreszcie, choć pewność
porządkuje pytania sensownie, jako bezwzględna miara jest źle
wykalibrowana: model jest systematycznie zbyt pewny siebie. Ta ostatnia
obserwacja godzi sprzeczne stanowiska z literatury - PPL jest
użytecznym sygnałem porządkowym, ale zawodnym sygnałem ilościowym.

Dalsza część pracy ma następującą strukturę. Sekcja~\ref{sec:tlo}
omawia tło teoretyczne i powiązane prace. Sekcja~\ref{sec:metodologia}
opisuje zbiór danych, modele, sposób skoringu i badane metryki.
Sekcja~\ref{sec:wyniki} przedstawia wyniki korelacyjne w trzech
ustawieniach, a sekcja~\ref{sec:kalibracja} - analizę kalibracji.
Sekcja~\ref{sec:dyskusja} zestawia obserwacje z literaturą, a
sekcja~\ref{sec:limitations} omawia ograniczenia.
